%% ---------------------------------------------------------------------
%%  FIGURA - Algoritmo de um passo de simulacao (DG-06)
%% ---------------------------------------------------------------------
\begin{figure}[htb]
	\caption{Algoritmo de um passo de simulação, com destaque para a ordem de ADR-004}
	\label{fig:tick}
	\centering
	\begin{tikzpicture}[
			font=\footnotesize,
			acao/.style={draw, rounded corners=2pt, minimum width=6.4cm, minimum height=0.78cm, align=center, fill=black!4},
			critico/.style={acao, fill=black!14, very thick},
			saida/.style={draw, rounded corners=2pt, minimum width=3.2cm, minimum height=0.78cm, align=center, fill=black!8},
			dec/.style={draw, diamond, aspect=3.0, inner sep=1pt, align=center, font=\scriptsize},
			>=Stealth,
			fl/.style={->, thick},
			rot/.style={font=\scriptsize, inner sep=2pt, fill=white}
		]

		\node[acao] (a1) at (0,0) {1. Aplicar a próxima direção válida da fila};
		\node[acao, below=0.42cm of a1] (a2) {2. Calcular a nova cabeça: cabeça $+$ direção};
		\node[dec, below=0.6cm of a2] (d1) {cabeça fora\\ da grade?};
		\node[acao, below=0.6cm of d1] (a3) {3. Verificar se a cabeça alcançou o alvo};
		\node[critico, below=0.42cm of a3] (a4) {4. Se não cresceu: remover a cauda};
		\node[dec, below=0.6cm of a4] (d2) {cabeça em célula\\ do corpo?};
		\node[acao, below=0.6cm of d2] (a5) {5. Inserir a nova cabeça no início do corpo};
		\node[acao, below=0.42cm of a5] (a6) {6. Se cresceu: somar pontos, reduzir o intervalo, sortear novo alvo};

		\node[saida, right=1.5cm of d1] (o1) {Encerrar partida};
		\node[saida, right=1.5cm of d2] (o2) {Encerrar partida};
		\node[saida, right=1.5cm of a6, align=center] (o3) {Ganhar\\ (sem célula livre)};

		\draw[fl] (a1) -- (a2);
		\draw[fl] (a2) -- (d1);
		\draw[fl] (d1) -- node[rot, left] {não} (a3);
		\draw[fl] (d1) -- node[rot, above] {sim} (o1);
		\draw[fl] (a3) -- (a4);
		\draw[fl] (a4) -- (d2);
		\draw[fl] (d2) -- node[rot, left] {não} (a5);
		\draw[fl] (d2) -- node[rot, above] {sim} (o2);
		\draw[fl] (a5) -- (a6);
		\draw[fl] (a6) -- (o3);
	\end{tikzpicture}
	\legend{Fonte: elaborado pelo autor (2026) a partir de \texttt{index.html}, seção~8.
	O passo 4 é destacado porque sua posição \emph{antes} do teste de colisão (ADR-004) determina se
	ocupar a célula recém-liberada pela cauda é sobrevivência ou morte.}
\end{figure}
